% =============================================================================
% CS224R Final Report Template
% =============================================================================
% This template is structured for your CS224R final report with two main parts:
%   1. One-page extended abstract (required as first page)
%   2. Main paper (recommended ~8 pages)
%
% The extended abstract summarizes your project's main findings and accomplishments.
% The main paper details your methods, prior work, contributions, and results.
%
% For project groups, include a breakdown of individual contributions.
% =============================================================================

\documentclass{article}

\usepackage[preprint]{cs224r_2025}
\usepackage[utf8]{inputenc} % allow utf-8 input
\usepackage[T1]{fontenc}    % use 8-bit T1 fonts
\usepackage{hyperref}       % hyperlinks
\usepackage{url}            % simple URL typesetting
\usepackage{booktabs}       % professional-quality tables
\usepackage{amsfonts}       % blackboard math symbols
\usepackage{nicefrac}       % compact symbols for 1/2, etc.
\usepackage{microtype}      % microtypography
\usepackage{xcolor}         % colors
\usepackage{amsmath}
\usepackage{graphicx}       % for including figures
\usepackage{lipsum}         % for placeholder text

\title{In-context Search: Efficiency Boost or Fundamentally New Capability?}

\author{%
  Angel Raychev \\
  Department of Computer Science\\
  Stanford University\\
  \texttt{angelray@stanford.edu} \\
  % examples of more authors
  \And
  Yalcin Tur \\
  Department of Computer Science\\
  Stanford University\\
  \texttt{yalcintr@stanford.edu} \\
  \And
  Mihajlo Stojkovic\\
  Department of Computer Science\\
  Stanford University\\
  \texttt{mstojkov@stanford.edu} \\
}

\begin{document}

% ============================================================================
% PART 1: ONE-PAGE EXTENDED ABSTRACT
% ============================================================================


% ============================================================================
% PART 2: MAIN PAPER
% ============================================================================
\maketitle


\section{Introduction}

Reasoning models, most notably OpenAI's O1 \cite{openai-o1-2412.16720}, represent a novel category of large language models (LLMs) characterized by the incorporation of in-context search during generation. This approach uniquely enables scaling inference-time computation to directly improve model performance. While classical search methods also leverage additional computational resources to enhance standard LLM performance, reasoning models employing in-context search typically demonstrate superior computational efficiency, achieving comparable or improved results with significantly reduced computational cost.

Traditional language models estimate the probability of producing an answer \( a \) given a question \( q \) as follows:
\begin{equation}
    \pi_{\text{data}}(a \mid q) \propto \int \pi_{\text{data}}(a \mid qs_1 \dots s_n) \prod_{i=1}^n \pi_{\text{data}}(s_i \mid q s_1\ldots s_{i-1})\, dS
\end{equation}
where \( s_1, s_2, \dots, s_n \) represent intermediate solution steps. In contrast, reasoning models using in-context search formulate generation differently \cite{xiang2025towards}:
\begin{equation}
    \pi_{\text{data}}(s_1 \dots s_n a\mid q) \propto
\int \pi_{\text{data}}(s_1 \dots s_n a\mid q z_1 \dots z_K)  \prod_{i=1}^K \pi_{\text{data}}(z_i \mid q z_1\ldots z_{i-1})\, dZ
\end{equation}
Here, \( z_1, z_2, \dots, z_K \) represent intermediate, model-generated reasoning steps leading to the final solution \( s_1, \dots, s_n \). Structured search methods impose explicit structure on these intermediate steps. In-context search models, however, do not enforce internal structure during training but instead guide these intermediate reasoning steps externally.

A significant open question raised by \cite{xiang2025towards} is whether reasoning models utilizing in-context search merely provide greater computational efficiency compared to structured search methods or if they indeed possess the capability to solve fundamentally new classes of problems that classical structured search inherently cannot, regardless of available compute. To answer this, analyzing scaling laws for these methods becomes essential.


\section{Related Work}

\paragraph{Classical search.}
Structured-reasoning with LLMs has been tackled via prompt-based tree search methods such as Tree-of-Thought \cite{yao2023tree} and Reasoning via Planning \cite{hao-etal-2023-reasoning}, which use a frozen language model as a value function and perform BFS or DFS. These methods rely on prompt engineering and self-verification by the generator, leading to brittleness, and their shallow search constrains them to simple tasks. By contrast, AlphaZero–style Monte Carlo Tree Search trains a dedicated policy network alongside a separate value head, enabling deep, scalable exploration without depending on the model’s own probabilities. Empirically, this critic-based MCTS outperforms prompt-based verifiers and shallow tree search. MCTS seems the best of structured search, and that’s why we use it following the TS-LLM methodology \cite{feng2023alphazero}.

\paragraph{In-context search.}
Parallel to explicit search, a new family of \emph{in-context} models has emerged (OpenAI O1 \cite{openai-o1-2412.16720}, DeepSeek R1 \cite{deepseek-r1-2501.12948}, xAI Grok-3 \cite{xai-grok-3-2025}, etc.).  Unlike classical approaches, these systems solve problems by iterating entire trajectories inside the prompt window. The first systematic recipe is Stream of Search (SoS) \cite{gandhi2024stream}, which trains on the \textit{full} search traversals, not merely the final answers, and then refines the model via Self-Taught Reasoner (STaR) \cite{zelikman2022star}. Our work follows SoS for the in-context branch while adopting an AlphaZero-inspired MCTS for the classical branch, enabling a controlled head-to-head comparison.


\section{Methods}

\paragraph{Classical search (MCTS).}
We adopt Monte Carlo Tree Search (MCTS), similar to AlphaZero \cite{feng2023alphazero}.  
Instead of using  roll-outs, we train a policy~$\pi_\theta$ and a value
function~$v_\phi$ in together.
For every question~$q$, MCTS produces a trajectories
$s_1s_2\!\ldots s_n a$ that ends with a parsable answer~$a$.
Positive trajectories (those whose answer is correct) fine-tune the policy, while all terminal
trajectories supervise the value network in a binary-classification objective:

\begin{equation}\label{policy_loss}
  \mathcal{L}_{\text{policy}}
  = \mathbb{E}_{q,s,a\sim\mathcal{D}^+}
    \bigl[-\log \pi_\theta(s_1\!\ldots s_n a \mid q)\bigr],
\end{equation}
\begin{equation}\label{value_loss}
  \mathcal{L}_{\text{value}}
  = \mathbb{E}_{q,s,a\sim\mathcal{D}}
    \Bigl[
      \ell_{\text{BCE}}\!\bigl[v_\phi(s_1\!\ldots s_n a \mid q),\, r(a\mid q)\bigr]
      + \sum_{t=1}^{n}
        \ell_{\text{BCE}}\!\bigl[v_\phi(s_1\!\ldots s_t \mid q),\, r(a\mid q)\bigr]
    \Bigr],
\end{equation}

where $\mathcal{D}$ is the distribution of all terminal traces and
$\mathcal{D}^+$ the subset of positives.
The binary reward $r(a\mid q)\!\in\!\{0,1\}$ is provided by an oracle grader.
Training alternates between
(i)~data collection with the current MCTS policy/value and
(ii)~parameter updates via~\eqref{policy_loss}–\eqref{value_loss}.

\paragraph{MCTS details.}
Each search cycle has four phases—selection, expansion, evaluation, back-propagation.

\textit{Selection.}  
At depth~$i$ we pick the next action (sentence-level)
\[
  s_i^\star
  = \arg\max_b
    U(s_1\!\ldots s_{i-1}s_i^{(b)}),\qquad
  U = Q + c_{\text{explore}}
      \frac{\sqrt{N(s_1\!\ldots s_{i-1})}}
           {N(s_1\!\ldots s_{i-1}s_i^{(b)})+1},
\]
where $Q(\cdot)$ is the running average return and
$N(\cdot)$ the visit count.
We continue until reaching a leaf or forming a complete answer
$s_1\!\ldots s_n a$.

\textit{Expansion \& evaluation.}  
From a leaf we sample $B$ children
\[
  s_i^{(1)},\dots,s_i^{(B)}
  \sim \pi_\theta(\cdot\mid q s_1\!\ldots s_{i-1}),
\]
evaluate each with $v_\phi$, initialise
$N\!=\!0$ and $Q\!=\!v_\phi$, and skip this step if the node is already
terminal.

\textit{Back-propagation.}  
Traversing back to the root, we update visit counts and
running means:
\[
  Q(s_1\!\ldots s_{j-1}) \gets
  Q(s_1\!\ldots s_{j-1}) +
  \frac{v_\phi(s_1\!\ldots s_{j-1}\mid q) - Q(s_1\!\ldots s_{j-1})}
       {N(s_1\!\ldots s_{j-1}) + 1},
  \quad j\le i.
\]
Search stops when the compute budget is exhausted.
During training we store only answer-yielding paths; at inference we output the most-visited branch based on the Q-values.

\paragraph{In-context search (SoS + STaR).}
After convergence we generate full tree traversals with the
final~$\pi_\theta,v_\phi$, flatten them, and train a fresh
in-context model~$\overline{\pi}_\psi$ from scratch, following
Stream of Search \cite{gandhi2024stream}.
Traversals of multiple tree sizes serve as supervised data,
and the model is also trained to emit the final answer.
We then apply STaR for an additional performance boost.

\section{Implementation Details}

The complete source code is available at
\url{https://github.com/RaychevAngel/classical_vs_contextual_search}.
All components were implemented from scratch.  
Although an open-source MCTS existed, its lack of clean parallelism
motivated a fresh implementation.
Similarly, the original Stream-of-Search (SoS) codebase was small but
not fully aligned with our workflow, so we re-implemented it to keep the
repository self-contained.

\paragraph{Model architecture.}
We employ Qwen-2.5 checkpoints throughout.
The policy $\pi_\theta$ and value $v_\phi$ share the 1.5 B
parameter base, whereas the in-context model
$\overline{\pi}_\psi$ starts from the 3 B variant.
Both approaches therefore use an equal 3 B total parameters drawn from
the same family, ensuring a fair comparison.

\paragraph{Task and dataset.}
Our benchmark generalises \emph{Game 24}:
given four integers $x_1,\dots,x_4\!\in\![1,12]$
and a target $t\!\in\![1,999]$, build $t$ exactly once using
$\{+,-,\times,\div\}$ and parentheses.
We generated 105 000 problems and split them into
99 000/3 000/3 000 for train/val/test.
To prime the policy, we hand-crafted solutions for the first 9 000
training instances, instructing the model to emit each step on a new
line and to format the final answer parse-ably.

% --- SAMPLES: do not modify ---
\noindent
\begin{minipage}[t]{0.32\textwidth}
\scriptsize\begin{verbatim}
q: Use 1, 4, 8, 11 to make 308.
s1: 8-1=7 Left: 4, 7, 11
s2: 4*7=28 Left: 11, 28
s3: 11*28=308 Left: 308
a: The  answer is: 4*(8-1)*11=308.
\end{verbatim}
\end{minipage}
\hfill
\begin{minipage}[t]{0.32\textwidth}
\scriptsize\begin{verbatim}
q: Use 1, 3, 8, 12 to make 32.
s1: 3-1=2 Left: 2, 8, 12
s2: 2*12=24 Left: 8, 24
s3: 8+24=32 Left: 32
a: The answer is: (3-1)*12+8=32.
\end{verbatim}
\end{minipage}
\hfill
\begin{minipage}[t]{0.32\textwidth}
\scriptsize\begin{verbatim}
q: Use 5, 7, 10, 11 to make 339.
s1: 5*7=35 Left: 10, 35, 11
s2: 10*35=350 Left: 11, 350
s3: 350-11-339 Left: 339
a: The answer is: 5*7*10-11=339.
\end{verbatim}
\end{minipage}
\newline
% --- end samples ---

\paragraph{Bootstrapping.}
The base policy $\pi_\theta^0$ is fine-tuned for one epoch on the
synthetic solutions (LR = $5\!\times\!10^{-5}$, batch = 256) to obtain
$\pi_\theta^1$.
Pre-training $v_\phi$ on synthetic negatives is
possible but brittle—enumerating the many ways a derivation can fail
often traps the network in local minima—so we skip this step.

\paragraph{Iterative MCTS training.}
At iteration $i\!\ge\!1$ we:
First, we sample 9 000 unseen training questions. Next, we grow one tree per question using a branch factor of 5, up to 25 expansions, and an exploration constant $c_{\text{explore}}=0.3$. Finally, we collect all terminal traces and use them as the next minibatch. vLLM serves both models; the policy uses temperature 1.0
and top-5 nucleus sampling, filtering duplicates.
For the value head we take a designated logit and apply a sigmoid.

Large trees raise diversity and help learning; with more time we would
explore branch factors 10–20.
Each round trains
$\pi_\theta^i$ (batch = 256) and $v_\phi^i$ (batch = 1024) for one epoch,
starting at LR $=2\!\times\!10^{-5}$ and decaying to
$5\!\times\!10^{-6}$ by iteration 6, where performance plateaus.
Intermediate-state labels are \emph{soft}: instead of the binary reward
$r(a\mid q)$ we use the average final reward reachable from that node,
reducing gradient noise without changing optima
(cf.\ Eq.~\ref{value_loss}).

\paragraph{In-context model (SoS).}
We freeze $\pi_\theta^6,v_\phi^6$, sample 12 000 fresh questions,
and record complete MCTS traversals.
States are logged in first-visit order, then flattened into a single
string; the Q-optimal solution and its answer (if correct) are appended.
Extra markers help the model navigate the tree.

% --- TREE SAMPLE: do not modify ---
\noindent
\begin{minipage}[t]{0.3\textwidth}
\scriptsize{
\begin{verbatim}
q: Use 3, 3, 5, 12 to make 51.
<START_THOUGHT>
z1: N1->Q | 3*12=36 Left: 3, 5, 36
z2: N2->Q | 12*5=60 Left: 3, 3, 60
z3: N3->Q | 5*12=60 Left: 3, 3, 60
z4: N4->N1 | 3*5=15 Left: 15, 36
z5: N5->N2 | 60-3=57 Left: 3, 57
z6: N6->N3 | 3+3=6 Left: 6, 60
z7: N7->N1 | 5*3=15 Left: 15, 36
z8: N8->N3 | 60-3=57 Left: 3, 57
\end{verbatim}}
\end{minipage}
\hfill
\begin{minipage}[t]{0.34\textwidth}
\scriptsize{
\begin{verbatim}
z9: N9->N2 | 3*3=9 Left: 9, 60
z10: N10->N4 | 36+15=51 Left: 51
z11: N11->N10->N4->N1->Q | 3*12+3*5= 51
z12: N12->N7 | 36+15=51 Left: 51
z13: N13->N12->N7->N1->Q | 3*12+5*3= 51
z14: N14->N3 | 3*3=9 Left: 9, 60
z15: N15->N5 | 57-3=54 Left: 54
z16: N16->N6 | 60-6=54 Left: 54
z17: N17->N9 | 60-9=51 Left: 51
z18: N18->N17->N9->N2->Q | 12*5-3*3= 51
\end{verbatim}}
\end{minipage}
\hfill
\begin{minipage}[t]{0.3\textwidth}
\scriptsize{
\begin{verbatim}
z19: N19->N8 | 57-3=54 Left: 54
z20: N20->N15->N5->N2->Q | 12*5-3-3= 54
z21: N21->N14 | 60-9=51 Left: 51
z22: N22->N21->N14->N3->Q | 5*12-3*3= 51
z23: N23->N19->N8->N3->Q | 5*12-3-3= 54
<END_THOUGHT>
<START_ANSWER>
a: 3*12+3*5=51
<END_ANSWER>
\end{verbatim}}
\end{minipage}
\newline
% --- end tree sample ---

Traversals are generated with branch factors
$\{1,2,3,5,8,13,21\}$.
A separate $\overline{\pi}_\psi$ is trained for each
factor (16 epochs, LR $=2\!\times\!10^{-5}$, batch = 64).

\paragraph{STaR refinement.}
We apply two STaR variants to the best
$\overline{\pi}_\psi$:
We compare two STaR variants on 12 000 new questions. In the first (STaR), we generate exactly one trajectory per question. In the second (STaR\*), we generate up to ten trajectories for each question and select the first one that produces a correct answer. Both converge in a single iteration, producing
$\overline{\pi}_\psi^{\text{STaR}}$ and
$\overline{\pi}_\psi^{\text{STaR}^\ast}$.


\section{Experiments and Results}

We retain five MCTS checkpoints (the initial policy-only run is omitted, as it lacks a value head) and
\(3\times 7\) SoS checkpoints—\(\overline{\pi}_\psi\),
\(\overline{\pi}_\psi^{\mathrm{STaR}}\), and
\(\overline{\pi}_\psi^{\mathrm{STaR}^\ast}\)—for every branch factor
\(b\in\{1,2,3,5,8,13,21\}\).
All models are evaluated on the held-out test set; results appear in
Figures~\ref{mcts} and~\ref{sos_mcts}.

\begin{figure}[h]
    \centering
    \begin{minipage}[t]{0.48\textwidth}
        \centering
        \includegraphics[width=\linewidth]{mcts.png}
        \caption{MCTS at different checkpoints}
        \label{mcts}
    \end{minipage}
    \hfill
    \begin{minipage}[t]{0.48\textwidth}
        \centering
        \includegraphics[width=\linewidth]{sos_comparison.png}
        \caption{Classical vs.\ In-Context Search}
        \label{sos_mcts}
    \end{minipage}
\end{figure}

MCTS get's up to \(51.13\%\) accuracy at its final checkpoint, improving steadily as the compute budget grows.The in-context search variants outperform MCTS at small and medium budgets, yet their accuracy drops abruptly beyond \(\sim20{,}000\) tokens—likely a context-length limitation—so MCTS ultimately scales better.

\section{Discussion and Conclusion}

Our study set out to clarify whether in-context search confers
\emph{intrinsically} greater reasoning power or merely a more efficient
use of compute relative to classical structured search.
The evidence so far is nuanced.

\paragraph{Compute–accuracy trade-off.}
Stream-of-Search (SoS) beats policy–value MCTS at very small inference
budgets, but the advantage vanishes once the budget rises:  
SoS plateaus, whereas MCTS keeps climbing and ultimately overtakes.
Hence, the apparent efficiency of in-context search does not translate
into better asymptotic performance under our current training regime.

\paragraph{Limitations of our SoS training.}
We applied only a single STaR pass—essentially a supervised warm-start.
The model was never \emph{reinforced} to avoid poor trajectories.
Incorporating policy-gradient updates (e.g.\ PPO, DPO, GRPO) or
contrastive penalties against dead-ends could push the in-context model
much further up the compute curve.

\paragraph{Task bias.}
Generalised \textsc{Game 24} favours classical search:
\begin{itemize}
\item a sparse, well-typed action space simplifies semantic filtering;
\item trees are shallow, so MCTS explores them exhaustively;
\item vector search over known formulas already yields strong priors.
\end{itemize}
Moving to broader mathematical domains—with larger vocabularies,
longer derivations, and looser syntax—should accentuate the
flexibility of in-context reasoning.

\paragraph{Context-length ceiling.}
The sharp drop in SoS accuracy beyond $\sim$20\,k prompt tokens almost
certainly stems from context-length limits.
Techniques such as retrieval-augmented prompting, segment-wise
compression, or hierarchical deliberation may alleviate this barrier.

\paragraph{Outlook.}
While MCTS remains the most scalable \emph{structured} search in our
experiments—supporting the TS-LLM philosophy—the question of whether
in-context search can ultimately outstrip external search is still open.
Future work should:
\begin{enumerate}
\item integrate stronger RL objectives into SoS training,
\item evaluate on tasks with deeper reasoning chains,
\item test longer-context or memory-augmented language models.
\end{enumerate}
Only with such advances can we definitively determine whether
in-context search unlocks qualitatively new capabilities or simply
reshuffles the compute–performance frontier.



\section{Team Contributions}
\begin{itemize}
  \item \textbf{Angel Raychev:} Lead developer. Implemented most of the codebase, repeatedly refactored the MCTS engine, and ran all experiments.
  \item \textbf{Yalcin Tur:} Contributed to the initial code base.
  \item \textbf{Mihajlo Stojković:} Contributed to the initial code base.
\end{itemize}

\noindent\textit{Note.}  
This work extends our CS224N project. All three authors built the first
MCTS prototype, after which Yalcin and Mihajlo stepped away.  
Angel has since rewritten the entire codebase several times, fixing
numerous bugs and spending hundreds of hours tuning MCTS
hyper-parameters.  
The Stream-of-Search module is completely new and introduced further
engineering challenges. Development is still ongoing.



\bibliographystyle{ACM-Reference-Format}
\bibliography{reference}

\end{document}